\section{Market Disruptions}

\begin{frame}{Supply Shocks}
    \mytext{
    The shape of the demand curve is pretty much a given, at least in the short run.

    However, the supply curve can be shifted by external shocks, such as war, taxes, subsidies...
    It can then
    \begin{itemize}
        \item shift up or down: every unit is equally affected (i.e. unit tax)
        \item pivot: marginal costs increase (ad-valorem tax)
    \end{itemize}
    }
\end{frame}

\begin{frame}{Supply Shocks: Ad valorem vs. per unit}
    \mycomment{
        The supply curve moves when the marginal costs of a unit change. There are generally two ways this can happen:
        \begin{itemize}
            \item the good becomes more expensive per unit, for example by a sudden increase in the price for a material that goes into it, or by tariff or tax levied per unit.
            \item there is a tax levied on the value of the goods sold, an \textit{ad valorem} tax. This pivots the supply curve. Import tariffs are often levied this way.
        \end{itemize}
        Whether external shock or tax policy: both effects can (per chance or per design) result in the same new equilibrium with the same new volume and price.

        This means that all welfare effects we will find for a shifting supply curve will hold for a pivoting supply curve, everything else being equal (ceteris paribus). 
    }
    \showmytikz[1,2,3]{../shared/tikz/tikz-supplyshiftorpivot}
    
\end{frame}

\begin{frame}{Welfare Effect of Supply Shocks}
    \mycomment{
    Consider the symmetric situation. Now every unit produced suddenly costs a little more. That means the supply curve shifts upward. This could for example be a tax levied.
    }
    \showmytikz[1]{../shared/tikz/tikz-demandandsupply-supplyshock}
\end{frame}

\begin{frame}{Welfare Effect of Supply Shocks}
    \mycomment{
    The result is a smaller welfare altogether. In the symmetric case, both consumers and procuers are worse off. You could say the damage has been shared. If it was a tax, consumers and producers share the burden of the tax (equally).
    }
    \showmytikz[2]{../shared/tikz/tikz-demandandsupply-supplyshock}
\end{frame}

\begin{frame}{Welfare Effect: inelastic Demand}
    \mycomment{Consider again the situation with highly  inelastic demand and elastic supply.

    }
    \showmytikz[1,2]{../shared/tikz/tikz-demandandsupplysimple-inelasticdemand}
\end{frame}

\begin{frame}{Welfare Effect: inelastic Demand}
    \mycomment{

    The biggest surplus is the consumer surplus to not surprise. Consumers are willing to pay much more, but the do not have to.    

    Here, a shift of the supply curve will affect the consumer more than the supplier. The less elastic party carries more of the burden. 
    }
    \showmytikz[1,3,4]{../shared/tikz/tikz-demandandsupplysimple-inelasticdemand}
\end{frame}

\begin{frame}{Welfare Effect: Elastic Demand}
    \mycomment{
    Consider again demand be elastic and supply inelastic, so that originally the bigger surplus goes to the producers. 
    }

    \showmytikz[1,2]{../shared/tikz/tikz-demandandsupplysimple-inelasticsupply}
\end{frame}

\begin{frame}{Welfare Effect: Elastic Demand}
    \mycomment{
    Again, an unit tax will be carried more by the inelastic party; in this case, the suppliers.
    }

    \showmytikz[1,3,4]{../shared/tikz/tikz-demandandsupplysimple-inelasticsupply}
\end{frame}

\begin{frame}{Rigid Quantity Restriction}
    \mycomment{
    The producer surplus will change. Some suppliers will be forced out of the market (for no apparent reason). Producers will gain surplus from the consumer, but also lose some surplus due to the smaller quantity. The size of the overall effect---if the producer surplus goes up or down---will again depend on the elasticities.
    }
    \showmytikz[1]{../shared/tikz/tikz-quantityrestriction}
\end{frame}

\begin{frame}{Rigid Quantity Restriction}
    \mycomment{

    The producer surplus will change. Some suppliers will be forced out of the market (for no apparent reason). Producers will gain surplus from the consumer, but also lose some surplus due to the smaller quantity. 
    
    The size of the overall effect - if the producer surplus goes up or down - will depend on the elasticities.

    Here, we have modelled the supply curve $SS_2$ perfectly inelastic. This is not realistic, as there will usually be some holders of the good willing to part at higher price. 
    }
    \showmytikz[2]{../shared/tikz/tikz-quantityrestriction}
\end{frame}

\begin{frame}{Summing Up: Welfare Effects of Shocks}
    \mytext{
    \begin{itemize}
        \item the new price and the new quantity depend on the elasticities of supply and demand
        \item neither suppliers nor buyers "decide to pass this on": the market decides
        \item a cost shock will likely never be passed on to 100\% to the market price
        \item the less elastic party bears most of the burden of a shock
        \item vice versa holds: all of the above can be applied for positive and negative shocks.
        
    \end{itemize}
    }
    \mycomment{
    Consider for example a tax that was levied on energy, which is highly inelastic. Such a tax has advantages for the state since consumers can hardly void it (there is no easy substitute). That means stable income from the tax, but also little distortion on the markets.

    If now the state decides to reduce this tax, the price for consumers will change almost in the same height. As consumers were bearing most of the burden of the tax before, the profite more of the tax reduction as well.
    }
\end{frame}
